% TOP_PROBLEM: {"problemId": "problem.long-range-order-in-the-quantum-heisenberg-ferromagnet", "canonicalTitle": "Long-Range Order in the Quantum Heisenberg Ferromagnet", "releaseRank": 202, "primaryDomain": "mathematical_physics", "primaryDomainLabel": "Mathematical physics", "displayStatus": "open_with_solved_subcases", "releaseStatus": "open_with_solved_subcases"}
% !TeX program = lualatex
% Definitions and sources retained from research_batch_201_202.tex; research is in GPT_6_Astra_Ultra.
\documentclass[11pt]{article}
\usepackage[margin=25mm]{geometry}
\usepackage{fontspec}
\setmainfont{Latin Modern Roman}
\usepackage{amsmath,amssymb,mathtools}
\usepackage{unicode-math}
\setmathfont{Latin Modern Math}
\usepackage{enumitem,needspace}
\usepackage{xurl,hyperref}
\hypersetup{colorlinks=true,urlcolor=blue}
\setcounter{secnumdepth}{1}
\setlength{\parindent}{0pt}
\setlength{\parskip}{5pt}
\setlength{\emergencystretch}{3em}
\setlist[itemize]{leftmargin=3em,itemsep=5pt,topsep=4pt,parsep=0pt}
\allowdisplaybreaks
\newcommand{\N}{\mathbb{N}}
\newcommand{\Z}{\mathbb{Z}}
\newcommand{\Q}{\mathbb{Q}}
\newcommand{\R}{\mathbb{R}}
\newcommand{\C}{\mathbb{C}}
\newcommand{\F}{\mathbb{F}}
\newcommand{\A}{\mathbb{A}}
\newcommand{\PP}{\mathbb P}
\newcommand{\E}{\mathbb{E}}
\newcommand{\eps}{\varepsilon}
\newcommand{\dd}{\,\mathrm d}
\newcommand{\sref}[2]{\hyperlink{source:#1:#2}{[S#2]}}
\newcommand{\eref}[2]{\hyperlink{extra:#1:#2}{[E#2]}}
% Turn the next switch off for a shorter reading copy. The quotations preserve
% every exactTarget field, including qualifications omitted from a short summary.
\newif\ifshowcatalogue
\showcataloguetrue
\newcommand{\cataloguescope}[1]{\ifshowcatalogue
	\paragraph{Exact scope in the supplied catalogue.}
	\begin{quote}\small #1\end{quote}\fi}
\begin{document}
\setcounter{section}{201}
	% ================================================================
	% problemId: problem.long-range-order-in-the-quantum-heisenberg-ferromagnet
	% source releaseRank: 202; source releaseStatus: open_with_solved_subcases
	% exactTarget SHA-256: 6b6155a73c536e1df98675d36c115216d36a32dc833614a37b7c949626ad07c8
	\clearpage
	\section{\texorpdfstring{Long-Range Order in the Quantum Heisenberg Ferromagnet}{Long-Range Order in the Quantum Heisenberg Ferromagnet}}\label{202}
	\subsection{Definitions and mathematical statement}
	At each site of $\Z^d$, spin-$S$ operators act on $\C^{2S+1}$, satisfy the angular-momentum commutation relations, and have $\mathbf S_x^2=S(S+1)I$. The finite-volume Hamiltonian is
	$H_\Lambda=-\sum_{\langle x,y\rangle\subset\Lambda}\mathbf S_x\cdot\mathbf S_y$.
	An infinite-volume equilibrium state is understood in the quantum Gibbs/KMS sense of the cited model. For each $d\ge3$ and positive half-integer $S$, the assertion is that some finite $\beta_0(d,S)$ satisfies
	\[
	\forall\beta\ge\beta_0\text{ finite},\quad
	\exists\text{ translation-invariant zero-field equilibrium state }\omega_\beta,
	\quad\omega_\beta(\mathbf S_0)\ne0.
	\]
	The state is allowed to break rotational symmetry; its rotation-invariant mixture need not have nonzero magnetization. The coupling is fixed to one.
	\par\smallskip\textit{Formulation and concept references:} \sref{202}{1}.
	\cataloguescope{For the nearest-neighbor isotropic spin-S quantum Heisenberg ferromagnet on Z\textasciicircum{}d with Hamiltonian H=-sum\_\{<x,y>\} S\_x dot S\_y, assert: for every integer d>=3 and every S in\{1/2,1,3/2,...\}, there is a finite beta\_0(d,S)>0 such that for each finite beta>=beta\_0 the infinite-volume zero-field model admits a translation-invariant equilibrium Gibbs state with nonzero spin expectation. Nonzero expectation refers to spontaneous magnetization in a symmetry-breaking state, not the rotation-invariant mixture. Coupling J=1 fixes temperature units; no pointwise-correlation equivalence is asserted.}
	\subsection{Short English statement}
	At sufficiently low positive temperature, must the quantum Heisenberg ferromagnet in every dimension at least three possess an equilibrium state with spontaneous magnetization?
	\subsection{Sources}
	\begin{itemize}\raggedright
		\item[\textnormal{[S1]}]\hypertarget{source:202:1}{} Long range order for the quantum Heisenberg model, contributed27January1999.
		\url{https://web.math.princeton.edu/~aizenman/OpenProblems_MathPhys/9901.HeisenbergFerr.html}.
		{\small\textit{Catalogue formulation source.}}
		% BEGIN RESEARCH REFERENCES 202
		\item[\textnormal{[E1]}]\hypertarget{extra:202:1}{}
		M. Correggi, A. Giuliani, and R. Seiringer, \emph{Validity of the
			spin-wave approximation for the free energy of the Heisenberg
			ferromagnet}, Communications in Mathematical Physics 339 (2015),
		279--307, Theorem 2.1, Section 3, and Proposition 5.2.
		\url{https://arxiv.org/pdf/1312.7873v2};
		\url{https://doi.org/10.1007/s00220-015-2402-0}.
		{\small\textit{Published result; arXiv version 2 is dated 10 August 2019.
				Consulted 24 September 2026.}}
		\item[\textnormal{[E2]}]\hypertarget{extra:202:2}{}
		E. R. Speer, \emph{Failure of reflection positivity in the quantum
			Heisenberg ferromagnet}, Letters in Mathematical Physics 10 (1985),
		41--47, abstract, Introduction, and the finite-volume counterexamples.
		\url{https://sites.math.rutgers.edu/~speer/homepage/pubs/LMP_10_41.pdf}.
		{\small\textit{Its specified regimes are not replaced here by a claim
				about every temperature or every possible reflection. Consulted
				24 September 2026.}}
		\item[\textnormal{[E3]}]\hypertarget{extra:202:3}{}
		D. Ueltschi, \emph{Random loop representations for quantum spin systems},
		Journal of Mathematical Physics 54 (2013), 083301; author manuscript
		arXiv:1301.0811, Theorem 3.2, the spin/loop correspondence of Section 3,
		and Section 7.1, especially equation (7.1).
		\url{https://arxiv.org/pdf/1301.0811}.
		{\small\textit{The spin-$1/2$ ferromagnetic representation is used, not
				the different models for which this paper proves long-range order.
				Consulted 24 September 2026.}}
		\needspace{9\baselineskip}
		\item[\textnormal{[E4]}]\hypertarget{extra:202:4}{}
		B. Nachtergaele and R. Sims, \emph{Lieb--Robinson bounds in quantum
			many-body physics}, in R. Sims and D. Ueltschi (eds.), \emph{Entropy
			and the Quantum}, Contemporary Mathematics 529, American Mathematical
		Society (2010), 141--176, Section 3.1, Theorem 3.1 and its proof
		(norm convergence of bounded-interaction dynamics on local observables).
		\url{https://arxiv.org/pdf/1004.2086v1}.
		{\small\textit{Locality input for the thermodynamic-limit argument;
				the Gaussian-smearing passage of the KMS identity is given in the text.
				Consulted 24 September 2026.}}
		\item[\textnormal{[E5]}]\hypertarget{extra:202:5}{}
		D. Heson, S. Starr, and J. Thornton, \emph{Violation of ferromagnetic
			ordering of energy levels in spin rings for the singlet},
		arXiv:2307.12773v5, 7 January 2025, abstract and Section 2.
		\url{https://arxiv.org/html/2307.12773v5}.
		{\small\textit{The manuscript distinguishes its largely numerical
				ordering claim from a rigorous singlet-sector result; neither is
				used as a proof of the lattice conjecture here. Consulted
				24 September 2026.}}
		% END RESEARCH REFERENCES 202
	\end{itemize}
\end{document}
